\scanpage{25}
Soit un $v\in\operatorname{End}_C(E_C)$
Il existe $u\in\operatorname{End}_{\deleted{\unsure{A}}B}(E)$
tel que $v+\deleted{\unsure{u\otimes_A B}}u_C$ soit \uline{inversible}.

\uline{Dém} Si $E\simeq B^{(N)}$, on a $\operatorname{End}_B(E)\simeq E^N$,
donc la remarque sur le relèvt d'éléments $v_i\in
\underline{\mathrm{End}}(E)(k(\mathfrak m_i))$ en un $v\in\underline{\mathrm{End}}(E)(B)$ reste valable
[on est ramené à relever des éléments \deleted{\unsure{des}}
$a_i\in E(k(\mathfrak m_i))$, et finalement, pr $E\simeq B^{(I)}$, des
éléments des $k(\mathfrak m_i)$ eux-$\overline{\mathrm m}\mathrm s$]. De $\overline{\mathrm m}$, grâce à
l'hypothèse que $\operatorname{rad}(C)$ nilpotent, il s'ensuit que
l'inversibilité ne dépend que de $w\bmod\operatorname{rad}(C)$.
Cela nous ramène au cas où $B$, \inserted{$C$} composés de corps,
\deleted{puis finalement}, \inserted{puis} grâce à l'hypothèse (ii), on a
alors $B\xrightarrow{\sim}C$, et on peut encore \unsure{$w|v|u$}$=\mathrm{id}_E-u$.
\deleted{\unsure{au où} $B\unsure{\subset|=|\rightrightarrows}C$}
\deleted{Dans le cas (i), on prend $u=\lambda\,\mathrm{id}_E$, \illegible{1}}

N.B. Il faudrait dégager la possibilité d'une
extension résiduelle \unsure{transcend.t|transcendante} $k(t)$, qui fait
le $\overline{\mathrm m}$ effet que (i) ---
% archive sheet number: 24
